\documentclass[main]{subfiles}
\usepackage[sols]{workbook}
\numbering{ws3-4}

\begin{document}

\section{Separable Differential Equations} \label{ws3-4}

\begin{framed}
\centerline{\textbf{Lesson Goals}}

You should be able to:
\begin{itemize}
\item Determine if a differential equation is separable
\item Find equilibrium solutions to a differential equation
\item Use the method of separation of variables to find the general solution to a separable differential equation


\end{itemize}
\end{framed}

  \begin{problemonly}
    \begin{framed}
    A first-order differential equation is called {\bf separable}
    if it can be written in the form $f(y) \frac{dy}{dt} = g(t)$.  
    \end{framed}
  \end{problemonly}  


\begin{enumerate}


  \item   \label{identify-separable}
  Decide whether each of the following differential equations is
  separable. 
  \begin{enumerate}      

    \item 
    \begin{problem}[0.5in]
      $\displaystyle \frac{dy}{dt} = e^{5t - 8y}$
    \end{problem}

    \begin{solution}
      Yes.  By laws of exponents, we can express the right-hand side as a product $e^{5t} e^{-8y}$. Thus the equation can be rewritten:
      \[e^{8y}\frac{dy}{dt} = e^{5t}\]
    \end{solution}

    \item
    \begin{problem}[0.5in]
    $\displaystyle \frac{dy}{dt} = t + y$.
    \end{problem}

    \begin{solution}
    No.  There is no way to express the right-hand side as a {\em product} 
    of the form $f(t) g(y)$.
    \end{solution}
    
    \item
    \begin{problem}[0.5in]
    $\displaystyle \frac{dy}{dt} = \frac{y^2 + 1}{\sin t}$.
    \end{problem}

    \begin{solution}
    Yes.  The right-hand side is $\displaystyle \underbrace{\frac{1}{\sin t}}_{f(t)} \: \underbrace{(y^2 + 1)}_{g(y)}$, so the differential equation becomes
    \[\frac{1}{y^2 + 1}\frac{dy}{dt} = \frac{1}{\sin t}\]
    \end{solution}
    
    \item
    \begin{problem}[1in]
    The logistic equation $\displaystyle \frac{dP}{dt} = \frac{1}{10} P \left( 1 - \frac{P}{800} \right)$.
    \end{problem}

    \begin{solution}
    Yes.  Any autonomous first-order DE such as this one is separable.
    \end{solution}
    
  \end{enumerate}

\pspace{\newpage}

\item \begin{problem}[2.5in]
  Find the general solution of the differential equation $\displaystyle
  \frac{dy}{dt} =ky$, where $k$ is a non-zero constant.
  \end{problem}

 \begin{solution} 
   One solution is the equilibrium solution $y(t) = 0$.  
   To find the other solutions, we use separation of variables:
   $$
     \frac{1}{y} \frac{dy}{dt} = k.
   $$
   Now integrate with respect to $t$:
   $$
     \int \frac{1}{y} \frac{dy}{dt}~dt = \int k~dt,
   $$
   or 
   $$
     \int \frac{1}{y}~dy = \int k~dt.
   $$
   Calculating the integrals, we have
   $$
     \ln |y| = kt + C,
   $$
   where $C$ is a constant of integration.  Exponentiate both sides
   to obtain
   $$
     |y| = e^{kt + C} = e^{C} e^{kt} = A e^{kt},
   $$
   where $A = e^{C} > 0$ is a constant.  It follows that
   $$
      y = \pm Ae^{kt} \mbox{ where } A > 0.
   $$
   Putting this family of solutions together with the equilibrium
   solution we found previously, the general solution of the 
   differential equation is given by $y = C_{1} e^{kt}$ where
   $C_{1}$ is {\em any} real constant.  (Note that $C_{1} = 0$
   is allowed, in which case we generate the equilibrium solution.)
 \end{solution}
 
 \item \begin{problem}[2.5in]
  Find the general solution of the differential equation $\displaystyle
  \frac{dy}{dt} =-y + 1$.
  \end{problem}

 \begin{solution} 

  We'll offer two ways to solve this differential equation.
  \begin{itemize}
    \item \textbf{Using separation of variables:}

    A good habit when using separation of variables is to first look for
    equilibrium solutions.  (It's easy to accidentally lose these
    solutions when separating; you won't lose anything if you're very
    careful not to divide by 0 anywhere, but it's easy to overlook
    things.)  Remember that equilibrium solutions are solutions where
    $M(t)$ is a constant function, so $\frac{dy}{dt} = 0$.  In this
    case, the only equilibrium solution is $y(t) = 1$.

    If $y \neq 1$, then $1 - y \neq 0$, and we can separate to
    get $\frac{1}{1 - y} \frac{dy}{dt} = 1$.\footnote{If we do this when
    $1 -  y = 0$, we are dividing by 0, which is bad.}  We can
    rewrite this as \[- 1 \cdot \frac{1}{y - 1} \frac{dy}{dt} =
    1.\footnote{You don't \textit{have} to simplify like this before
    integrating, but it might make integrating easier!}\] Integrating
    both sides with respect to $t$, we have \[- \ln | y - 1| = t + C.\] 
    Multiplying
    both sides by $-1$ gives \[\ln | y - 1| = - t -  C.\] Since
    $-C$ is just an arbitrary constant, we can give it a new name;
    let's call it $A$.  So,
    \begin{align*}
    \ln|y - 1| & = A - t \\
    e^{\ln |y - 1|} & = e^{A -  t} \\
    |y - 1| & = e^A e^{-t} \\
    y - 1 & = \pm e^{A} e^{- t}
    \end{align*} Again, $\pm e^A$ is just an arbitrary constant, so
    let's call it $B$.\footnote{Since $B = \pm e^A$, $B$ should
    technically be non-zero.}  So, \[y - 1 = B e^{- t},\] and
    \fbox{$y(t) = 1 + B e^{- t}$}.\footnote{The way we solved this,
    $B$ had to be non-zero.  However, at the very beginning, we
    pointed out that $y(t) = 1$ is a solution.  This solution could
    be written as $y(t) = 1 + B e^{- t}$ with $B = 0$.  So, in the
    final answer $y(t) = 1 + B e^{-t}$, $B$ really \textit{can} be
    any constant.}

    \item \textbf{Using a clever substitution:}

    We can rewrite the given differential equation as $\frac{dy}{dt} =
    - (y - 1)$, so let's substitute $N = y - 1$; then, $y = N +
    1$.  Since $y$ is a function of $t$, $N$ is also a function of $t$,
    and $\frac{dy}{dt} = \frac{dN}{dt}$.\footnote{This comes from taking
    $\frac{d}{dt}$ of both sides of the equation $y = N + 1$.  Or,
    you can think of it like this: $y(t)$ is just $N(t)$ plus a
    constant, so $y(t)$ and $N(t)$ must have the same derivative.}
    So, the original differential equation
    \begin{align*}
    \frac{dy}{dt} &= - (y - 1) \\
    \intertext{can be rewritten as}
    \frac{dN}{dt} &= - N
    \end{align*}
    Now, we can simply recognize that the solution of this differential
    equation is $N(t) = C e^{- t}$, because we've seen this sort of
    equation often.  Since $y = N + 1$, the general
    solution for the original differential equation is $y(t) = 1 + C
    e^{- t}$.
    
  \end{itemize}

  \end{solution}


  \item 
  \note{This can be done by substituting $N = M - 12$ or by doing
  separation of variables.  To motivate the particular choice of
  substitution, I look at the slope field; the slopes at $M = 12$ are
  all 0, so we might have an easier time if we moved the whole picture
  down by 12.}
  
  \begin{problem}[3in]
  Find the general solution of the differential equation $\displaystyle
  \frac{dM}{dt} = 2.4 - 0.2M$.
  \end{problem}

  \begin{solution} 
  Again, we'll offer two ways to solve this differential equation.
  \begin{itemize}
    \item \textbf{Using separation of variables:}

    A good habit when using separation of variables is to first look for
    equilibrium solutions.  (It's easy to accidentally lose these
    solutions when separating; you won't lose anything if you're very
    careful not to divide by 0 anywhere, but it's easy to overlook
    things.)  Remember that equilibrium solutions are solutions where
    $M(t)$ is a constant function, so $\frac{dM}{dt} = 0$.  In this
    case, the only equilibrium solution is $M(t) = 12$.

    If $M \neq 12$, then $2.4 - 0.2 M \neq 0$, and we can separate to
    get $\frac{1}{2.4 - 0.2 M} \frac{dM}{dt} = 1$.\footnote{If we do this when
    $2.4 - 0.2 M = 0$, we are dividing by 0, which is bad.}  We can
    rewrite this as \[- \frac{1}{0.2} \cdot \frac{1}{M - 12} \frac{dM}{dt} =
    1.\footnote{You don't \textit{have} to simplify like this before
    integrating, but it makes integrating easier!}\] Integrating
    both sides with respect to $t$, we have 
    \[- \frac{1}{0.2} \ln | M - 12| = t + C.\] Multiplying
    both sides by $-0.2$ gives \[\ln | M - 12| = -0.2 t - 0.2 C.\] Since
    $-0.2 C$ is just an arbitrary constant, we can give it a new name;
    let's call it $A$.  So,
    \begin{align*}
    \ln|M - 12| & = A - 0.2 t \\
    e^{\ln |M - 12|} & = e^{A - 0.2 t} \\
    |M - 12| & = e^A e^{-0.2 t} \\
    M - 12 & = \pm e^{A} e^{-0.2 t}
    \end{align*} Again, $\pm e^A$ is just an arbitrary constant, so
    let's call it $B$.\footnote{Since $B = \pm e^A$, $B$ should
    technically be non-zero.}  So, \[M - 12 = B e^{-0.2 t},\] and
    \fbox{$M(t) = 12 + B e^{-0.2 t}$}.\footnote{The way we solved this,
    $B$ had to be non-zero.  However, at the very beginning, we
    pointed out that $M(t) = 12$ is a solution.  This solution could
    be written as $M(t) = 12 + B e^{-0.2 t}$ with $B = 0$.  So, in the
    final answer $M(t) = 12 + B e^{-0.2t}$, $B$ really \textit{can} be
    any constant.}

    \item \textbf{Using a clever substitution:}

    We can rewrite the given differential equation as $\frac{dM}{dt} =
    -0.2 (M - 12)$, so let's substitute $N = M - 12$; then, $M = N +
    12$.  Since $M$ is a function of $t$, $N$ is also a function of $t$,
    and $\frac{dM}{dt} = \frac{dN}{dt}$.\footnote{This comes from taking
    $\frac{d}{dt}$ of both sides of the equation $M = N + 12$.  Or,
    you can think of it like this: $M(t)$ is just $N(t)$ plus a
    constant, so $M(t)$ and $N(t)$ must have the same derivative.}
    So, the original differential equation
    \begin{align*}
    \frac{dM}{dt} &= -0.2 (M - 12) \\
    \intertext{can be rewritten as}
    \frac{dN}{dt} &= -0.2 N
    \end{align*}
    Now, we can simply recognize that the solution of this differential
    equation is $N(t) = C e^{-0.2 t}$, because we've seen this sort of
    equation often.  Since $M = N + 12$, the general
    solution for the original differential equation is $M(t) = 12 + C
    e^{-0.2 t}$.
    
  \end{itemize}

  \end{solution}


%%%%%%%%%%%%%%%%%%%%%%%%%%%%%%%%%%%%%%%%%%%%%%%%%%%%%%%%%%%%%%%%%%%%%%  
  \item 
%  \note{This is a straightforward separation of variables.  For homework, I'll ask them to redo this question but with a different initial condition $y(0) = 1$ that gives a solution that blows up in finite time.}
  
  \begin{problem}[2.25in]
  Find the general solution of the
  differential equation $\displaystyle \frac{dy}{dt} = y^2$
  using separation of variables.  Then find the particular solution that
  satisfies the initial condition $y(0) = -1$.
  \end{problem}

  \begin{solution}
  There is one equilibrium solution:  $y = 0$.  Separation of
  variables helps us find the non-equilibrium solutions.  If $y \neq 0$,
  we may write
  $$
    \frac{1}{y^2} \frac{dy}{dt} \; = \; 1.
  $$
  Integrating both sides with respect to $t$,
  $$
    \int \frac{1}{y^2}~dy \; = \; \int 1~dt.
  $$
  Therefore, 
  $$
    - \frac{1}{y} \; = \; t + C,
  $$
  where $C$ is an arbitrary constant.  Thus, the general solution
  consists of the functions
  $$
    y = - \frac{1}{t + C} \qquad \mbox{and} \qquad y = 0.
  $$

  Using the initial condition $y(0) = -1$ yields $C = 1$, which means
  that the function satisfying both the differential equation and the
  initial condition is  
  $$
    y = \frac{-1}{t + 1}.
  $$
  \end{solution}

%%%%%%%%%%%%%%%%%%%%%%%%%%%%%%%%%%%%%%%%%%%%%%%%%%%%%%%%%%%%%%%%%%%%%%%%%

  \item 

  \note{This is a straightforward separation of variables, as long as they
  rewrite it as $y' = -2x y^2$.  This is an example where they're likely
  to lose the equilibrium solution $y = 0$.}
  
  \begin{enumerate}

  \item
  \begin{problem}[2in]
  Solve the differential equation $y' + 2x y^2 = 0$.  (Find all
  solutions.)
  \end{problem}

  \begin{solution}
  Let's rewrite the differential equation as $\frac{dy}{dx} = - 2x y^2$.
  This is separable, but let's first find the equilibrium solutions,
  where $y(x)$ is a constant and $\frac{dy}{dx} = 0$.  The only one is
  $y = 0$.

  If $y \neq 0$, we can separate variables to get $- y^{-2} \frac{dy}{dx} = 2x$.
  Integrating both sides, $y^{-1} = x^2 + C$, so $y = \frac{1}{x^2 +
  C}$.  So, the solutions are \fbox{$y = 0$ and $\displaystyle y =
  \frac{1}{x^2 + C}$}.
  \end{solution}

  \item
  \begin{problem}[1in]
  Find the solution of $y'+2x y^2=0$ that satisfies the initial condition $y(0)=0$. (Compare your answer to your answer to part (a).)     
  \end{problem} 

  \begin{solution}
  This is just the equilibrium solution $y(x) = 0$.  It's important not to miss this solution when you're solving for the general solution in part (a)!
  \end{solution}


  \item
  \begin{problem}[1in]
  Find the solution of $y'+2x y^2=0$ that satisfies the initial condition $y(0)=5$.
  \end{problem} 

  \begin{solution}
  From the previous part, we know the general solution of the differential equation is $y(x) = \frac{1}{x^2+C}$.  The initial value condition gives us that $5 = y(0) = \frac{1}{0^2+C} = \frac{1}{C}$, and therefore we must have $C=\frac{1}{5}$.  Therefore the particular solution to this initial value problem is $y = \frac{1}{x^2+\frac{1}{5}}$.
  \end{solution}

  \end{enumerate}

\pspace{\newpage}


  
  \item   

%  \note{I'll have the students do this problem.  They'll see another example of an ill-posed problem for homework (the leaky bucket equation backwards in time).}

  The main purpose for using differential equations as models of
  natural phenomena is to make predictions beyond what is experimentally
  measurable.  It's desirable that our models be {\em deterministic}:
  that is, we'd like there to be one and only one possible solution that could
  satisfy the differential equation and its initial conditions.  Without
  the luxury of determinism, a model's predictive power is questionable
  at best.  Sadly, not all models are deterministic, and here's an example.

  \begin{enumerate}
  \item
  \begin{problem}[0.75in]
    Consider the differential equation $\displaystyle \frac{dy}{dt} = 3 y^{2/3}$
    together with the initial condition $y(0) = 0$.  Can you think of a {\em really
    boring} example of a function that satisfies both the differential
    equation and its initial condition?    
  \end{problem}  

  \begin{solution}
    The constant function $y = 0$ is about as boring as it gets.
  \end{solution}

  \item
  \begin{problem}[2.5in]
    Use separation of variables to come up with a less boring example
    of a solution to the same differential equation subject to the same
    initial condition.  
  \end{problem}  

  \begin{solution}
    Separating the variables,
    $$
	\frac{1}{y^{2/3}} \frac{dy}{dt} = 3.
    $$
    Integrating both sides with respect to $t$, we obtain
    $$
	3y^{1/3} = 3t + C,
    $$
    where $C$ is a constant of integration.  From the initial condition,
    we see that $C = 0$.  Dividing both sides by $3$ and then cubing both
    sides, we obtain $y = t^3$.  It's easy to check that $y = t^3$ is,
    in fact, a solution of the differential equation and satisfies the
    initial condition!

    {\bf Discussion:}  Contrast this with the solution produced in Part (a)
    of the problem.  Since $y = 0$ and $y = t^3$ both qualify as solutions,
    this problem is {\em non-deterministic}.  It would be useless to try
    to make predictions of future behavior for a problem such as this one.    

  \end{solution}
  \end{enumerate}

  \item 
  \note{This is straightforward to separate, but students may not quite
  remember how to do the integration.}
  
  \begin{problem}[2.5in]
  Solve the differential equation $\displaystyle \frac{dy}{dt} =
  \frac{t^2}{y^2 (1 + t^3)}$.
  \end{problem}

  \begin{solution}
  This is separable.  We first look for equilibrium solutions.  These
  are solutions where $\frac{dy}{dt} = 0$ and $y$ is a constant.  So,
  these solutions must satisfy $0 = \frac{t^2}{y^2 (1 + t^3)}$, but
  there is no constant value of $y$ which makes this happen.  Therefore,
  there are no equilibrium solutions.

  Then, we separate: $y^2 \frac{dy}{dt} = \frac{t^2}{1 + t^3}$.  Integrating
  both sides with respect to $t$, we have 
  $\displaystyle \frac{1}{3} y^3 = \int \frac{t^2}{1 +
  t^3}~dt$.  To do the integral, let's substitute $u = 1 + t^3$.
  Then, $du = 3t^2~dt$, so
  \begin{align*}
  \frac{1}{3} y^3 &= \int \frac{1}{3} \cdot \frac{1}{u}~du \\
  &= \frac{1}{3} \ln |u| + C \\
  &= \frac{1}{3} \ln \left|1 + t^3\right| + C \\
  \shortintertext{Multiplying both sides by $3$,}
  y^3 &= \ln \left|1 + t^3\right| + 3C \\
  \shortintertext{Cube rooting both sides,}
  y &= \boxed{\left( \ln \left|1 + t^3\right| + 3C \right)^{1/3}}
  \end{align*}
  Of course, $3C$ is an arbitrary constant, so it's fine to give it a
  new name like $A$. 
  \end{solution}


\end{enumerate}




\end{document}


\begin{problemonly}
\newpage
\end{problemonly}

  \item 
  \note{A good thing to discuss is what
  should happen in the long run; that is, what do they think
  $\displaystyle \lim_{t \to \infty} S(t)$ is?}

  \begin{problem}[2in]
  You are mixing up a big batch
  of eggnog for the holidays.  You have a big mixing vat
  containing 100 gallons of egg yolks.  At time $t = 0$, you
  start pouring a milk and sugar mixture containing 1/2 lb of sugar per
  gallon into the vat at a rate of 4 gallons per minute.  At the same
  time, mixed eggnog starts leaving the vat a rate of 4 gallons per
  minute.

  Write a differential equation for the amount $S(t)$ of sugar (measured
  in pounds) in the vat at time $t$.  Can you write down an initial
  condition as well?
  \end{problem}

  \begin{solution}
  See the solutions to {Problem Set 29}, {\#1}.
  \end{solution}

  \item 
  \begin{problem}[1.5in]
  Suppose instead that, in the previous problem, mixed eggnog left the
  vat at a rate of 6 gallons per minute.  How would this change the
  differential equation and initial condition?
  \end{problem}

  \begin{solution}
  The basic framework is still the same:
  \[ (\textrm{rate of change of sugar in vat}) = (\textrm{rate at which
  sugar enters vat}) - (\textrm{rate at which sugar leaves vat}).\]
  The rate at which sugar enters the vat is exactly the same as before,
  so it is still 2 lb per minute.  The rate at which sugar leaves the
  vat can still be calculated as
  \begin{align*}
  {} & (\textrm{rate at which sugar is leaving vat at time $t$}) \\
  = {} & (\textrm{rate at which eggnog is leaving vat at time $t$})
  \cdot (\textrm{concentration of sugar in eggnog at time $t$}) \\
  = {} & 6 \cdot \left( \frac{\textrm{amount of sugar in
  vat at time $t$}}{\textrm{amount of eggnog in vat at time
  $t$}} \right) \\ 
  = {} & 6 \cdot \left( \frac{S}{\textrm{amount of eggnog in vat at
  time $t$}}
  \right)
  \end{align*}
  Since stuff is being added to the vat at 4 gallons per minute, but
  stuff is leaving the vat at 6 gallons per minute, the amount of eggnog
  in the vat is being reduced by 2 gallons per minute.  Thus, after $t$
  minutes, there are $100 - 2t$ gallons left in the vat.  Therefore, the
  rate at which sugar is leaving the vat at time $t$ is $6 \left(
  \frac{S}{100 - 2t} \right)$.  So, we have the final differential
  equation $\boxed{\frac{dS}{dt} = 2 - 6 \frac{S}{100 - 2t}}$.  The
  initial condition is still the same, $\boxed{S(0) = 0}$ since there is
  initially no sugar in the vat.
  \end{solution}





  % \item Note: the next problem introduces a method for solving non-separable differential equations using a clever substitution.  You aren't expected to learn this method nor to be able to do it on the exams.  We won't cover it in class.  But we're including this problem on the worksheet in case some of you want to see the method, and because you might want to look back at its solution to help solve one of the optional problems later this week. 

  % \label{first-hard-substitution-to-make-separable}
  % (Completely Optional Material) In this problem, we'll solve the non-separable differential equation $y'
  % + y = t e^{-t}$ by using a substitution $u = e^t y$.  
  % \begin{enumerate}
  %   \item
  %   \begin{problem}[0.25in]
  %   You are given a substitution $u(t) = e^t y(t)$; rewrite this
  %   equation to solve for $y(t)$.  
  %   \end{problem}

  %   \begin{solution}
  %   We just multiply both sides of the equation $u(t) = e^t y(t)$ by
  %   $e^{-t}$ to get $\boxed{y(t) = e^{-t} u(t)}$.
  %   \end{solution}
    
  %   \item
  %   \note{The main thing to look out for is students treating $u$ as a
  %   constant rather than a function of $t$ when differentiating.}
    
  %   \begin{problem}[1.25in]
  %   Use your answer to {{\#4}(a)} to rewrite the
  %   differential equation $y' + y = t e^{-t}$ in terms of $u$ (and $t$).
  %   \end{problem}

  %   \begin{solution}
  %   We found in {{\#4}(a)} that $y = e^{-t} u$.
  %   Let's plug this into the differential equation everywhere we see
  %   $y$.

  %   Differentiating $y = e^{-t} u$ (and remembering to use the Product
  %   Rule and Chain Rule), $y' = e^{-t} u' - e^{-t} u$.  So, the equation
  %   \begin{align*}
  %   y' + y &= t e^{-t} \\
  %   \shortintertext{can be rewritten as}
  %   (e^{-t} u' - e^{-t} u) + e^{-t} u &= t e^{-t} \\
  %   \shortintertext{Simplifying,}
  %   e^{-t} u' &= t e^{-t} \\
  %   \shortintertext{Multiplying both sides by $e^t$,}
  %   u' &= t
  %   \end{align*}
  %   So, our final differential equation is just $u' = t$. 
  %   \end{solution}

  %   \item

  %   \note{Students may want to use separation of variables here.  This
  %   certainly works, but it's nice to point out that it's not really
  %   necessary..}
    
  %   \begin{problem}[0.5in]
  %   Solve the differential equation you found in
  %   {{\#4}(b)}.  (Find all solutions.)
  %   \end{problem}

  %   \begin{solution}
  %   Our differential equation was $u' = t$, or $\frac{du}{dt} = t$.  You
  %   can solve by separating, but it's easier if you just realize that
  %   this equation is saying, ``$u$ is a function of $t$ whose derivative
  %   is the function $t$.''  In other words, $u$ is antiderivative of
  %   $t$, so $\boxed{u = \frac{1}{2} t^2 + C}$.
  %   \end{solution}

  %   \item
  %   \begin{problem}[0.25in]      
  %   Use your answer to the previous part to solve for $y(t)$.
  %   \end{problem}

  %   \begin{solution}
  %   We know from {{\#4}(a)} that $y(t) = e^{-t}
  %   u(t)$.  Since we found that $u = \frac{1}{2} t^2 + C$, this says
  %   that $\boxed{y(t) = \frac{1}{2} t^2 e^{-t} + C e^{-t}}$.
  %   \end{solution}

  %   \item
  %   \begin{problem}[1in]
  %   Find the solution of $y' + y = t e^{-t}$ satisfying the initial
  %   condition $y(0) = -1$. 
  %   \end{problem}

  %   \begin{solution}
  %   We know the general solution of $y' + y = te^{-t}$ is $y(t) =
  %   \frac{1}{2} t^2 e^{-t} + C e^{-t}$.  We can use the initial
  %   condition to solve for $C$.  If we plug in $t = 0$ and use the fact
  %   that $y(0) = -1$, we get $-1 = \frac{1}{2} 0^2 e^{-0} + C e^{-0} =
  %   C$, so $C = -1$.  Therefore, the solution we are looking for is
  %   $\boxed{y(t) = \frac{1}{2} t^2 e^{-t} - e^{-t}}$.

  %   \textit{Note:} You can always check that you have the right answer
  %   by plugging back into the original differential equation and initial
  %   condition.  Sometimes, this involves a lot of messy computation, so
  %   we've provided an applet
  %   \isitehref{http://isites.harvard.edu/fs/docs/icb.topic843879.files/docs/applets/SlopeFieldAndFunctionGrapher.nb}{\texttt{SlopeFieldAndFunctionGrapher.nb}}\footnote{If
  %   this opens in your web browser, go to File $\to$ Save As, save it
  %   to your computer, and then open it in Mathematica.} that makes it
  %   easy for you to visually check your answer.  Here's a screenshot of
  %   the applet:
  %   \begin{center}
  %   \includegraphics[scale=0.5]{images/diffeq-sub1-ivp-correct}
  %   \end{center}
  %   The applet graphs a slope field together with a function.  In this
  %   case, our differential equation is $y' + y = t e^{-t}$, or $y' = t
  %   e^{-t} - y$.  So, in the first box (highlighted in yellow), we enter
  %   $t e^{-t} - y$, which we type as \verb!t * E^(-t) - y!.  The
  %   function we want to graph is the solution we've found, $\frac{1}{2}
  %   t^2 e^{-t} - e^{-t}$; we enter this in the second box (highlighted
  %   in blue) as \verb!1/2 * t^2 * E^(-t) - E^(-t)!.  We can see that the
  %   function does indeed follow the slope field, and it passes through
  %   the point $(0, -1)$, which was our initial condition.  Therefore, it
  %   appears that we have found the right solution.

  %   For illustration purposes, let's imagine that we had done the
  %   problem completely incorrectly and ended up with a weird answer like
  %   $\arctan \left( \sqrt{1 + \cos t} \right)$.  If we instead entered
  %   this in the applet (as \verb!ArcTan[Sqrt[1 + Cos[t]]]!), we'd get a
  %   picture like this:
  %   \begin{center}
  %   \includegraphics[scale=0.5]{images/diffeq-sub1-ivp-incorrect}
  %   \end{center}
  %   We can see right away that this isn't right: the slope field should
  %   be tangent to the solution curve, but it isn't at all.  (Look, for
  %   instance, near $(-3, 0)$.)  
  %   \end{solution}
  % \end{enumerate}



  % \item \label{second-hard-substitution-to-make-separable}
  % \begin{problem}[0.5in]
  % Solve the differential equation $\displaystyle y' = \frac{y^2 +
  % xy}{x^2}$ by using the substitution $\displaystyle u = \frac{y}{x}$.
  % \end{problem}

  % \begin{solution}
  % Since $u = \frac{y}{x}$, $y = xu$.  Differentiating both sides (and
  % using the Product Rule on the right side), $y' = x u' + u$.  So, the
  % differential equation
  % \begin{align*}
  % y' &= \frac{y^2 + xy}{x^2} \\
  % \shortintertext{can be rewritten as}
  % xu' + u &= \frac{(xu)^2 + x (xu)}{x^2} \\
  % &= \frac{x^2 u^2 + x^2 u}{x^2} \\
  % &= u^2 + u
  % \shortintertext{Subtracting $u$ from both sides,}
  % xu' &= u^2
  % \shortintertext{Dividing both sides by $x$,}
  % u' &= \frac{u^2}{x}
  % \end{align*}
  % This is now a separable differential equation $\frac{du}{dx} =
  % \frac{u^2}{x}$.  As usual, we first look for equilibrium solutions:
  % the only one is $u(x) = 0$.  If $u \neq 0$, we can separate variables
  % to get $u^{-2}~du = \frac{1}{x}~dx$.  Integrating both sides, $-u^{-1}
  % = \ln |x| + C$, so $u = - \frac{1}{\ln |x| + C}$.

  % Since $y = xu$, the solutions of the differential equation
  % $\displaystyle y' = \frac{y^2 + xy}{x^2}$ are \fbox{$y = 0$ and
  % $\displaystyle y = - \frac{x}{\ln |x| + C}$}.
  % \end{solution}



%%%%%%%%%%%%%%%%%%%%%%%%%%%%%%%%%%%%%%%%%%%%%%%%%%%%%%%%%%%%%%%%%%%%%%  
  \item 
  \note{This is a straightforward separation of variables; it's fine for
  students to say the solutions are $y = \pm \sqrt{A - t^2}$, although it's
  worth asking (but not dwelling on) the question of what happens for
  large $t$.}
  
  \begin{problem}[2.25in]
  Solve the differential equation $\displaystyle \frac{dy}{dt} = -
  \frac{t}{y}$ using separation of variables.  (Previously, we solved by
  drawing the slope field, guessing the solution, and checking that it
  worked by plugging into the differential equation.)
  \end{problem}

  \begin{solution}
  There are no equilibrium solutions here (there is no constant value of
  $y$ which makes $- \frac{t}{y} = 0$ for all $t$).  We can separate
  variables to get \[y \frac{dy}{dt} = - t.\] Integrating both
  sides with respect to $t$, we have 
  \[\frac{1}{2} y^2 = - \frac{1}{2} t^2 + C.\] Multiplying both
  sides by $2$, \[y^2 = -t^2 + 2C.\] Since $2C$ is still just an
  arbitrary constant, we can give it a new name; let's call it $A$.
  So, \[y^2 = A - t^2,\] and the solutions are 
  \[ y = \pm \sqrt{A - t^{2}}.\]  Note that the solution curves
  follow semi-circles.  Interestingly, the solutions fail to exist
  once $t$ gets too large.
  \end{solution}

